%!TEX root = ./../../springer-mono-format/book.tex
%!TEX root = ./../../thesis-trier.tex




\clearpage

\IN{idempotent transformation}
We call $t \colon X \to X$ an \emph{idempotent (transformation)}
if it is idempotent w.r.t.~functional composition, i.e., we have $t^2 = t$.
For example, the idempotent transformations of defect one are precisely the
simple idempotents introduced in \cref{subsec:tail}.

First, in~\cref{prelim:prop:idempotent}, we list several equivalent conditions
to idempotency, and \cref{prelim:cor:idempotent} collects some consequences.
Then \cref{prelim:prop:compositionidempt} characterizes idempotency of the
composition, and
\cref{prelim:cor:compositionidempt-one,prelim:cor:comp-idempotents-idempotent,prelim:cor:comp-idempotents-idempotent2}
give characterizations under additional assumptions on the idempotency of the
given functions.
\Cref{prelim:trans:idemp:prop:inv-dupl} gives a formula for the inverse image
of the duplicate set of the composition involving at least one idempotent,
stating that it equals a union of inverse images of the individual
transformations.

\begin{proposition}\label{prelim:prop:idempotent}
	Let $t \colon X \to X$. Then the following are equivalent.
	\begin{enumerate}
		\item \label{prop:idemp:item:idemp} $t$ is idempotent.
		
		\item \label{prop:idemp:item:inv} $t(x) \in t^{-1}(t(x))$ for all $x \in X$.
		
		\item \label{prop:idemp:item:invy} $t^{-1}(y) \ne \emptyset$ implies $y \in t^{-1}(y)$ for $y \in X$.
		
		\item \label{prop:idemp:item:img} $t(X) \subseteq \fix{t}$.
		
		\item \label{prop:idemp:item:excl} $\move{t} \subseteq \excl{t}$.
		
		\item \label{prop:idemp:item:tmove}  $t(\move{t}) \subseteq \fix{t}$.
	\end{enumerate}
\end{proposition}
\begin{proof}
	The equivalences
	$\ref{prop:idemp:item:idemp} \Leftrightarrow~\ref{prop:idemp:item:inv}
	\Leftrightarrow~\ref{prop:idemp:item:invy}\Leftrightarrow~\ref{prop:idemp:item:img}$
	follow readily.
	
	$\ref{prop:idemp:item:img} \Rightarrow~\ref{prop:idemp:item:excl}$:
	This follows as $\move{t} = X \setminus \fix{t} \subseteq X \setminus t(X) = \excl{t}$.
	
	$\ref{prop:idemp:item:excl} \Rightarrow~\ref{prop:idemp:item:tmove}$:
	Here $t(\move{t}) \subseteq t(X) = X \setminus \excl{t} \subseteq X \setminus \move{t} = \fix{t}$.
	
	$\ref{prop:idemp:item:tmove} \Rightarrow~\ref{prop:idemp:item:img}$:
	As $t(\fix{t}) = \fix{t}$ and $t(\move{t}) \subseteq \fix{t}$, we find
	$t(X) = t(\fix{t}) \cup t(\move{t}) \subseteq \fix{t}$.
\end{proof}

\begin{corollary} \label{prelim:cor:idempotent}
	Suppose $t \colon X \to X$ is idempotent. Then we have the following:
	\begin{enumerate}
		\item \label{prelim:idemp:2} $t(X) = \fix{t}$,
		
		\item \label{prelim:idemp:3} $\move{t} = \excl{t}$,
		
		\item \label{prelim:idemp:1} $\dupl{t} = t(\move{t})$
		(equivalently, $t(x) \notin \dupl{t}$ implies $t(x) = x$),
		
		\item \label{prelim:idemp:4b} $\dupl{t} = t(\dupl{t})$, i.e., $t$ fixes every element of $\dupl{t}$,
		
		\item \label{prelim:idemp:5} $t^{-1}(\dupl{t}) = \move{t} \uplus \dupl{t}$,
		
		\item \label{prelim:idemp:7} $t(A) \subseteq A$ for $A \subseteq X$ with $\dupl{t} \subseteq A$.
	\end{enumerate}
\end{corollary}
\begin{proof}
	\ref{prelim:idemp:2}: This follows from \cref{prop:idemp:item:img} of
	\cref{prelim:prop:idempotent} and \cref{prelim:eqn:fiximage}.
	
	\ref{prelim:idemp:3}: This follows from \cref{prop:idemp:item:excl} of
	\cref{prelim:prop:idempotent} and \cref{prelim:eqn:exclmove}.
	
	\ref{prelim:idemp:1}: By \cref{prop:idemp:item:tmove} of
	\cref{prelim:prop:idempotent} and \cref{prelim:eqn:fixmovedupl} we have the
	following inclusions, which give the claim:
	\[
	t(\move{t}) \subseteq \fix{t}\cap t(\move{t}) \subseteq \dupl{t} \subseteq t(\move{t}).
	\]
	
	\ref{prelim:idemp:4b}: This follows as
	$\dupl{t} \subseteq t(X) \subseteq \fix{t}$ by \cref{prop:idemp:item:img} of
	\cref{prelim:prop:idempotent}.
	
	\ref{prelim:idemp:5}: By \cref{prelim:idemp:1} we have
	$\move{t} \subseteq t^{-1}(\dupl{t})$, and $\dupl{t} \subseteq t^{-1}(\dupl{t})$
	by \cref{prelim:idemp:4b}.
	Conversely, let $x \in t^{-1}(\dupl{t})$ with $x \notin \move{t}$.
	Then $x = t(x) \in \dupl{t}$.
	The union is disjoint as $\dupl{t} \subseteq \fix{t}$.
	
	\ref{prelim:idemp:7}: For $x \in A$ we have $t(x) = x \in A$ or, by
	\cref{prelim:idemp:1}, $t(x) \in \dupl{t} \subseteq A$.
\end{proof}

Next, we investigate when the composition of two transformations gives an
idempotent transformation.

\begin{proposition} \label{prelim:prop:compositionidempt}
	Let $t, h \colon X \to X$. Then the following are equivalent.
	\begin{enumerate}
		\item \label{prelim:prop:compositionidempt:1} $\concat{t}{h}$ is idempotent.
		
		\item \label{prelim:prop:compositionidempt:3}
		$\move{t} \cap h(t(X)) \subseteq \fix{\concat{t}{h}}$ and $\fix{t} \cap h(t(X)) \subseteq \fix{h}$.
		
		\item \label{prelim:prop:compositionidempt:2} $t(y) \in h^{-1}(y)$ for all $y \in h(t(X))$.
		
		\item \label{prelim:prop:compositionidempt:4} $h(\move{h} \cap t(X)) \subseteq \fix{\concat{t}{h}}$
		and $\fix{h} \cap t(X) \subseteq \fix{\concat{t}{h}}$.
	\end{enumerate}
\end{proposition}
\begin{proof}
	$\ref{prelim:prop:compositionidempt:1} \Rightarrow \ref{prelim:prop:compositionidempt:3}$:
	By \cref{prelim:prop:idempotent} we have
	$h(t(X)) \subseteq \fix{\concat{t}{h}}$, which gives the first inclusion;
	the second follows with \cref{prelim:trans:eqn:fix-comp}.
	
	$\ref{prelim:prop:compositionidempt:3} \Rightarrow \ref{prelim:prop:compositionidempt:2}$:
	Let $y \in h(t(X))$.
	If $y \in \move{t}$, then $h(t(y)) = y$ by the first assumption.
	Otherwise $y \in \fix{t}$, so $h(y) = y$ by the second assumption and hence
	$h(t(y)) = h(y) = y$.
	
	$\ref{prelim:prop:compositionidempt:2} \Rightarrow \ref{prelim:prop:compositionidempt:4}$:
	Let $x \in X$. Applying \ref{prelim:prop:compositionidempt:2} to
	$y = h(t(x))$ gives $h(t(x)) \in \fix{\concat{t}{h}}$, and applying it to
	$y = t(x)$ in case $t(x) \in h(t(X))$ gives $t(x) \in \fix{\concat{t}{h}}$.
	The first covers $h(\move{h} \cap t(X))$, the second $\fix{h} \cap t(X)$, as
	$t(x) \in \fix{h}$ implies $t(x) = h(t(x)) \in h(t(X))$.
	
	$\ref{prelim:prop:compositionidempt:4} \Rightarrow \ref{prelim:prop:compositionidempt:1}$:
	With $t(X) = ( \move{h} \cap t(X) ) \uplus ( \fix{h} \cap t(X) )$ we get
	\[
	h(t(X)) = h(\move{h} \cap t(X)) \cup ( \fix{h}\cap t(X) ) \subseteq \fix{\concat{t}{h}},
	\]
	so $\concat{t}{h}$ is idempotent by \cref{prelim:prop:idempotent}.
\end{proof}

\begin{corollary} \label{prelim:cor:compositionidempt-one}
	\label{prelim:prop:compositionidempth} \label{prelim:prop:compositionidemptt}
	Let $t, h \colon X \to X$.
	\begin{enumerate}
		\item \label{prelim:prop:compositionidempth:1} For idempotent~$h$,
		$\concat{t}{h}$ is idempotent if and only if
		$\move{t} \cap h(t(X)) \subseteq \fix{\concat{t}{h}}$.
		
		\item \label{prelim:prop:compositionidemptt:1} For idempotent~$t$,
		$\concat{t}{h}$ is idempotent if and only if
		$h(\move{h} \cap t(X)) \subseteq \fix{\concat{t}{h}}$.
	\end{enumerate}
\end{corollary}
\begin{proof}
	\ref{prelim:prop:compositionidempth:1}: For idempotent~$h$ we have
	$h(t(X)) \subseteq h(X) \subseteq \fix{h}$ by \cref{prelim:prop:idempotent},
	so the second condition of \cref{prelim:prop:compositionidempt:3} of
	\cref{prelim:prop:compositionidempt} always holds.
	
	\ref{prelim:prop:compositionidemptt:1}: For idempotent~$t$ we have
	$t(X) \subseteq \fix{t}$ by \cref{prelim:prop:idempotent}, so
	$\fix{h} \cap t(X) \subseteq \fix{h} \cap \fix{t} \subseteq \fix{\concat{t}{h}}$
	by \cref{prelim:trans:eqn:fix-comp}, and the second condition of
	\cref{prelim:prop:compositionidempt:4} of
	\cref{prelim:prop:compositionidempt} always holds.
\end{proof}

\begin{remark}
	By \cref{prelim:prop:compositionidempth:1} of
	\cref{prelim:cor:compositionidempt-one}, for idempotent~$h$ the composition
	$\concat{t}{h}$ is idempotent if $h(t(X)) \subseteq \fix{t}$, as then
	$\move{t} \cap h(t(X)) = \emptyset$.
	This also follows directly from
	$h(t(X)) \subseteq \fix{h} \cap \fix{t} \subseteq \fix{\concat{t}{h}}$ with
	\cref{prop:idemp:item:img} of \cref{prelim:prop:idempotent}.
	The condition is not necessary: if $h(x) = c$ for all $x \in X$ and
	$t(c) \ne c$, then $h$ and $\concat{t}{h}$ are idempotent and
	$h(t(X)) \cap \fix{t} = \emptyset$.
\end{remark}

\begin{corollary} \label{prelim:cor:comp-idempotents-idempotent}
	Let $t, h \colon X \to X$ be idempotent.
	Then $\concat{t}{h}$ is idempotent if and only if
	$\move{t} \cap h(\move{h} \cap t(X)) \subseteq \fix{\concat{t}{h}}$.\footnote{Equivalently,
		for $x \in X$, $h(t(x)) \ne t(x)$ and $t(h(t(x))) \ne h(t(x))$ imply
		$h(t(x)) = h(t(h(t(x))))$.}
\end{corollary}
\begin{proof}
	If $\concat{t}{h}$ is idempotent, then
	$h(t(X)) \subseteq \fix{\concat{t}{h}}$ by \cref{prelim:prop:idempotent},
	and the set in question is contained in $h(t(X))$.
	
	Conversely, as~$t$ is idempotent we have
	$\fix{h}\cap t(X) = \fix{h}\cap \fix{t} \subseteq \fix{\concat{t}{h}}$
	by \cref{prelim:trans:eqn:fix-comp}, and as~$h$ is idempotent we have
	$\fix{t} \cap h(t(X)) \subseteq \fix{t}\cap \fix{h} \subseteq \fix{\concat{t}{h}}$.
	With $t(X) = ( \move{h} \cap t(X) ) \uplus ( \fix{h} \cap t(X) )$ and
	$X = \move{t} \uplus \fix{t}$ we obtain
	\begin{multline*}
		h(t(X)) = ( \move{t} \cap h(\move{h} \cap t(X)) ) \cup
		\\ ( \move{t} \cap ( \fix{h}\cap t(X) ) ) \cup ( \fix{t} \cap h(t(X)) ),
	\end{multline*}
	where we used $h(\fix{h}\cap t(X)) = \fix{h}\cap t(X)$.
	All three sets are contained in $\fix{\concat{t}{h}}$, the first by
	assumption, so $\concat{t}{h}$ is idempotent by \cref{prelim:prop:idempotent}.
\end{proof}

\begin{corollary} \label{prelim:cor:comp-idempotents-idempotent2}
	Let $t, h \colon X \to X$ be idempotent.
	Then $\concat{t}{h}$ is idempotent if one of the following is true:
	\begin{enumerate}
		\item \label{prelim:cor:comp-idempotents-idempotent2:1}
		$\dupl{h} \subseteq \fix{\concat{t}{h}}$, i.e., $h(t(y)) = y$ for all $y \in \dupl{h}$,
		
		\item \label{prelim:cor:comp-idempotents-idempotent2:2}
		$\excl{t} \subseteq \fix{\concat{t}{h}}$,
		
		\item \label{prelim:cor:comp-idempotents-idempotent2:3} $ht = th$.
	\end{enumerate}
\end{corollary}
\begin{proof}
	\ref{prelim:cor:comp-idempotents-idempotent2:1}: As~$h$ is idempotent, we
	have $h(\move{h} \cap t(X)) \subseteq h(\move{h}) = \dupl{h}$ by
	\cref{prelim:idemp:1} of \cref{prelim:cor:idempotent}, which is contained in
	$\fix{\concat{t}{h}}$ by assumption.
	Thus the claim follows from \cref{prelim:prop:compositionidemptt:1} of
	\cref{prelim:cor:compositionidempt-one}.
	
	\ref{prelim:cor:comp-idempotents-idempotent2:2}: With
	\cref{prelim:prop:idempotent} we find
	$\move{t} = \excl{t} \subseteq \fix{\concat{t}{h}}$, so
	\cref{prelim:cor:comp-idempotents-idempotent} gives the claim.
	
	\ref{prelim:cor:comp-idempotents-idempotent2:3}: We have
	$h(t(h(t(x)))) = h(h(t(t(x)))) = h(t(x))$ for $x \in X$.
	Note that here only the idempotency of~$h$ and~$t$ in the last step is used,
	but not \cref{prelim:cor:comp-idempotents-idempotent}.
\end{proof}

\begin{proposition} \label{prelim:trans:idemp:prop:tX-in-fixh}
	Let $t, h \colon X \to X$ be idempotent.
	Then the following are equivalent.
	\begin{enumerate}
		\item \label{prelim:trans:idemp:prop:tX-in-fixh:1} $\concat{t}{h} = t$.
		
		\item \label{prelim:trans:idemp:prop:tX-in-fixh:2} $t(X) \subseteq \fix{h}$.
		
		\item \label{prelim:trans:idemp:prop:tX-in-fixh:3} $t(X) \cap h^{-1}(\dupl{h}) \subseteq \dupl{h}$.
	\end{enumerate}
\end{proposition}
\begin{proof}
	$\ref{prelim:trans:idemp:prop:tX-in-fixh:1} \Rightarrow \ref{prelim:trans:idemp:prop:tX-in-fixh:2}$:
	By assumption we have $h(t(x)) = t(x)$ for all $x \in X$.
	
	$\ref{prelim:trans:idemp:prop:tX-in-fixh:2} \Rightarrow \ref{prelim:trans:idemp:prop:tX-in-fixh:3}$:
	For $t(x) \in h^{-1}(\dupl{h})$ we have $t(x) = h(t(x)) \in \dupl{h}$.
	
	$\ref{prelim:trans:idemp:prop:tX-in-fixh:3} \Rightarrow \ref{prelim:trans:idemp:prop:tX-in-fixh:1}$:
	Suppose $h(t(x)) \ne t(x)$ for some $x \in X$.
	Then $h(t(x)) \in \dupl{h}$ by \cref{prelim:idemp:1} of
	\cref{prelim:cor:idempotent}, so $t(x) \in t(X) \cap h^{-1}(\dupl{h})$, but
	$t(x) \notin \dupl{h}$ as~$h$ fixes every element of $\dupl{h}$ by
	\cref{prelim:idemp:4b} of \cref{prelim:cor:idempotent}.
	Thus \cref{prelim:trans:idemp:prop:tX-in-fixh:3} fails.
\end{proof}

Next, we refine \cref{prelim:func:prop:inv-dupl} for the case that one
transformation is idempotent.
For idempotent $t \colon X \to X$ and $A \subseteq X$ we have
\begin{equation} \label{prelim:trans:idemp:eqn:inv}
	t^{-1}(A) = ( A \cap t(X) ) \cup t^{-1}( A \cap \dupl{t} ),
\end{equation}
as $x \in t^{-1}(A)$ with $x \notin \dupl{t}$ satisfies $x = t(x) \in A \cap t(X)$
by \cref{prelim:idemp:1} of \cref{prelim:cor:idempotent}, while
$A \cap t(X) \subseteq t^{-1}(A)$ holds as~$t$ fixes $t(X)$.

\begin{proposition}  \label{prelim:trans:idemp:prop:inv-dupl}
	Let $t \colon X \to X$ be idempotent and $h \colon X \to Y$. Then
	\begin{align*}
		(\concat{t}{h})^{-1}(\dupl{\concat{t}{h}})
		& = t^{-1}(\dupl{t}) \cup ( h^{-1}( \dupl{h_{|t(X)}} ) \cap t(X) ) \\
		& = t^{-1}(\dupl{t}) \cup h^{-1}( \dupl{\concat{t}{h}}).
	\end{align*}
\end{proposition}
\begin{proof}
	Set $A = h^{-1}(\dupl{h_{|t(X)}})$ and $B = A \cap t(X) = h_{|t(X)}^{-1}(\dupl{h_{|t(X)}})$.
	For the first identity we have
	\begin{align*}
		(\concat{t}{h})^{-1}(\dupl{\concat{t}{h}})
		& = t^{-1}(\dupl{t}) \cup t^{-1}(A) \tag{\cref{prelim:func:prop:inv-dupl}} \\
		& = t^{-1}(\dupl{t}) \cup B \cup t^{-1}(A \cap \dupl{t} ) \tag{Eqn.~\eqref{prelim:trans:idemp:eqn:inv}} \\
		& = t^{-1}(\dupl{t}) \cup B,
	\end{align*}
	the last step as $t^{-1}(A \cap \dupl{t}) \subseteq t^{-1}(\dupl{t})$.
	
	For the second identity, note that
	$\excl{t} \cup \dupl{t} \subseteq t^{-1}(\dupl{t})$ by
	\cref{prelim:idemp:3,prelim:idemp:5} of \cref{prelim:cor:idempotent}.
	Hence
	\begin{align*}
		t^{-1}(\dupl{t}) \cup h^{-1}(\dupl{\concat{t}{h}})
		& = t^{-1}(\dupl{t}) \cup h_{|t(X)}^{-1}(\dupl{\concat{t}{h}}) \tag{$\excl{t} \subseteq t^{-1}(\dupl{t})$} \\
		& = t^{-1}(\dupl{t}) \cup h_{|t(X)}^{-1}(h_{|t(X)}(\dupl{t})) \cup B \tag{Eqn.~\eqref{prelim:eqn:dupl0}} \\
		& = t^{-1}(\dupl{t}) \cup B,
	\end{align*}
	where the last step holds since, by \cref{prelim:eqn:finvfinclusion},
	$h_{|t(X)}^{-1}(h_{|t(X)}(\dupl{t})) \subseteq \dupl{t} \cup B$
	and $\dupl{t} \subseteq t^{-1}(\dupl{t})$.
	Together with the first identity this shows the claim.
\end{proof}

%\begin{proof}
%	Since $X \setminus t(X) = \excl{t} = \move{t} \subseteq t^{-1}(\dupl{t})$
%	by \cref{prelim:idemp:3,prelim:idemp:1} of \cref{prelim:cor:idempotent},
%	we have
%	\begin{equation} \label{prelim:trans:idemp:eqn:absorb}
%		t^{-1}(\dupl{t}) \cup S = t^{-1}(\dupl{t}) \cup ( S \cap t(X) )
%		\qquad \text{for every } S \subseteq X.
%	\end{equation}
%	Set $A = h^{-1}(\dupl{h_{|t(X)}})$ and $B = A \cap t(X)$.
%	
%	For the first identity, \cref{prelim:func:prop:inv-dupl} gives
%	$(\concat{t}{h})^{-1}(\dupl{\concat{t}{h}}) = t^{-1}(\dupl{t}) \cup t^{-1}(A)$.
%	By \cref{prelim:trans:idemp:eqn:absorb} we may replace $t^{-1}(A)$ by
%	$t^{-1}(A) \cap t(X)$, which equals~$B$ as~$t$ fixes $t(X)$.
%	
%	For the second identity, \cref{prelim:trans:idemp:eqn:absorb} likewise
%	replaces $h^{-1}(\dupl{\concat{t}{h}})$ by
%	$h^{-1}(\dupl{\concat{t}{h}}) \cap t(X)$, so it suffices to show that this
%	set is contained in $t^{-1}(\dupl{t}) \cup B$; the reverse inclusion holds
%	as $\dupl{h_{|t(X)}} \subseteq \dupl{\concat{t}{h}}$ by
%	\cref{prelim:eqn:dupl0}.
%	So let $x \in t(X)$ with $h(x) \in \dupl{\concat{t}{h}}$.
%	By \cref{prelim:eqn:dupl0} we have $h(x) \in \dupl{h_{|t(X)}}$, giving
%	$x \in B$, or $h(x) = h(y)$ for some $y \in \dupl{t}$.
%	In the latter case $y \in \dupl{t} \subseteq t(X)$, so either $x = y \in t^{-1}(\dupl{t})$,
%	or $x$ and~$y$ are distinct elements of $t(X)$ with the same image
%	under~$h$, whence $h(x) \in \dupl{h_{|t(X)}}$ and again $x \in B$.
%\end{proof}


\begin{remark}
	In \cref{prelim:trans:idemp:prop:inv-dupl}, the set
	$h^{-1}(\dupl{h_{|t(X)}}) \cap t(X)$ may be replaced with either
	$h^{-1}(\dupl{h_{|t(X)}})$ or $h^{-1}(\dupl{\concat{t}{h}}) \cap t(X)$,
	since both contain $h^{-1}(\dupl{h_{|t(X)}}) \cap t(X)$ and are contained in
	$h^{-1}(\dupl{\concat{t}{h}})$.
	However, $h^{-1}(\dupl{h_{|t(X)}})$ cannot be replaced with
	$h^{-1}(\dupl{h})$, see \cref{prelim:func:ex:inv-dupl}.
\end{remark}


\todo[inline]{auskommentiert}
\begin{comment}
\IN{idempotent transformation}
We call $t \colon X \to X$ an \emph{idempotent (transformation)}
if it is idempotent w.r.t.~functional composition, i.e., we have $t^2 = t$. 
For example, the idempotent transformations of defect one are precisely the simple idempotents introduced in \cref{subsec:tail}.

First, in~\cref{prelim:prop:idempotent}, we list several equivalent conditions to idempotency,
and \cref{prelim:cor:idempotent} collects some consequences.
Then \cref{prelim:prop:compositionidempt} characterizes idempotency
of the composition,
and \cref{prelim:prop:compositionidempth,prelim:prop:compositionidemptt,prelim:cor:comp-idempotents-idempotent2}
give characterizations under the additional assumptions on the idempotency of the given functions.
\Cref{prelim:trans:idemp:prop:inv-dupl} gives
a formula for the inverse image of the duplicate set of the composition involving at least one idempotent,
stating that it equals a union of inverse images of the individual transformations.

\begin{proposition}\label{prelim:prop:idempotent}
	Let $t \colon X \to X$. Then the following are equivalent.
	\begin{enumerate}
		\item \label{prop:idemp:item:idemp} $t$ is idempotent.
		
		\item \label{prop:idemp:item:inv} $t(x) \in t^{-1}(t(x))$ for all $x \in X$.
		
		\item \label{prop:idemp:item:invy} $t^{-1}(y) \ne \emptyset$ implies $y \in t^{-1}(y)$ for $y \in X$.
		
		\item \label{prop:idemp:item:img} $t(X) \subseteq \fix{t}$.
									
		\item \label{prop:idemp:item:excl} $\move{t} \subseteq \excl{t}$. 		

		\item \label{prop:idemp:item:tmove}  $t(\move{t}) \subseteq \fix{t}$.			
	\end{enumerate}
\end{proposition}
\begin{proof}
The equivalences $\ref{prop:idemp:item:idemp} \Leftrightarrow~\ref{prop:idemp:item:inv} 
\Leftrightarrow~\ref{prop:idemp:item:invy}\Leftrightarrow~\ref{prop:idemp:item:img}$ follow readily.
		
$\ref{prop:idemp:item:img} \Rightarrow~\ref{prop:idemp:item:excl}$: 
This follows as $\move{t} = X \setminus \fix{t} \subseteq X \setminus t(X) =  \excl{t}$.

$\ref{prop:idemp:item:excl} \Rightarrow~\ref{prop:idemp:item:tmove}$:
Here $t(\move{t}) \subseteq t(X) = X \setminus \excl{t} \subseteq X \setminus \move{t} = \fix{t}$.

$\ref{prop:idemp:item:tmove} \Rightarrow~\ref{prop:idemp:item:img}$: As $t(\fix{t}) = \fix{t}$
and $t(\move{t}) \subseteq \fix{t}$,
we find $t(X) = t(\fix{t} \cup \move{t}) = t(\fix{t}) \cup t(\move{t}) \subseteq \fix{t}$.
\end{proof} 

\begin{corollary} \label{prelim:cor:idempotent}
Suppose $t \colon X \to X$ is idempotent.
Then we have the following:
\begin{enumerate}
\item \label{prelim:idemp:2} $t(X) = \fix{t}$

\item \label{prelim:idemp:3} $\move{t} = \excl{t}$

\item \label{prelim:idemp:1} $\dupl{t} = t(\move{t})$ (equivalently, $t(x) \notin \dupl{t}$ implies $t(x) = x$)

\item \label{prelim:idemp:4b} $\dupl{t} = t(\dupl{t})$

\item \label{prelim:idemp:5} $t^{-1}(\dupl{t}) = \move{t} \cup \dupl{t}$

\item \label{prelim:idemp:7} For $A \subseteq X$ with $\dupl{t} \subseteq A$ we have $t(A) \subseteq A$.
\end{enumerate}	
\end{corollary}
\begin{proof}
\ref{prelim:idemp:2}: This follows from \cref{prop:idemp:item:img} of \cref{prelim:prop:idempotent} 
and \cref{prelim:eqn:fiximage}.
	
\ref{prelim:idemp:3}: This follows from \cref{prop:idemp:item:excl} of \cref{prelim:prop:idempotent} 
and \cref{prelim:eqn:exclmove}.

\ref{prelim:idemp:1}: By \cref{prop:idemp:item:tmove} of \cref{prelim:prop:idempotent} and \cref{prelim:eqn:fixmovedupl} we have the following inclusion which give the claim:
\[
t(\move{t})
\subseteq \fix{t}\cap t(\move{t})
\subseteq \dupl{t} 
\subseteq t(\move{t}).
\]

\ref{prelim:idemp:4b}: This follows as $\dupl{t} \subseteq t(X) \subseteq \fix{t}$
by \cref{prop:idemp:item:img} of \cref{prelim:prop:idempotent}.

\ref{prelim:idemp:5}: By \cref{prelim:prop:compinvAt} 
and since $\move{t} \subseteq t^{-1}(\dupl{t})$
and $\dupl{t} \subseteq t(X) \subseteq \fix{t}$ by \cref{prelim:idemp:1,prelim:idemp:2}, respectively,
we have 
\[
	t^{-1}(\dupl{t}) = ( \move{t} \cap t^{-1}(\dupl{t}) ) \cup ( \fix{t} \cap \dupl{t} ) 
	= \move{t} \cup \dupl{t}.
\]

\ref{prelim:idemp:7}: This follows from \cref{prelim:idemp:1}.
\end{proof}

Next, we investigate when the composition of two transformations gives an idempotent transformation.

\begin{proposition} \label{prelim:prop:compositionidempt}
Let $t, h \colon X \to X$. Then the following are equivalent.
	
\begin{enumerate} 
\item \label{prelim:prop:compositionidempt:1} $\concat{t}{h}$ is idempotent.

\item \label{prelim:prop:compositionidempt:3} 
$\move{t} \cap h(t(X)) \subseteq \fix{ht}$ and $\fix{t}  \cap h(t(X)) \subseteq \fix{h}$.

\item \label{prelim:prop:compositionidempt:2} $t(y) \in h^{-1}(y)$ for all $y \in h(t(X))$.

\item \label{prelim:prop:compositionidempt:4}  $h(\move{h} \cap t(X)) \subseteq \fix{\concat{t}{h}}$
and $\fix{h} \cap t(X) \subseteq \fix{\concat{t}{h}}$.
\end{enumerate} 
\end{proposition}
\begin{proof}
$\ref{prelim:prop:compositionidempt:1} \Rightarrow \ref{prelim:prop:compositionidempt:3}$: Let $\concat{t}{h}$ be idempotent.
Since, by \cref{prelim:prop:idempotent}, $h(t(X)) \subseteq \fix{\concat{t}{h}}$, we get 
$\move{t} \cap h(t(X)) \subseteq \fix{\concat{t}{h}}$.
With \cref{prelim:trans:eqn:fix-comp} 
we find $\fix{t} \cap h(t(X)) \subseteq \fix{h}$.

$\ref{prelim:prop:compositionidempt:3} \Rightarrow \ref{prelim:prop:compositionidempt:2}$: Let $x \in X$. Set $y = h(t(x))$
and consider $h(t(y))$. If $t(y) \ne y$, then $h(t(y)) = y$
by assumption. Otherwise, if $t(y) = y$, then $h(y) = y$
by assumption, so that $h(t(y)) = y$.
Hence, in each case, we find $t(y) \in h^{-1}(y)$.

$\ref{prelim:prop:compositionidempt:2} \Rightarrow \ref{prelim:prop:compositionidempt:4}$:
Let $x \in X$. Suppose $h(t(x)) \ne t(x)$.
By Item~\ref{prelim:prop:compositionidempt:2} applied to $y = h(t(x))$, we have $h(t(y)) = y$
and so $h(t(x)) \in \fix{\concat{t}{h}}$.
Next, suppose $h(t(x)) = t(x)$.
By Item~\ref{prelim:prop:compositionidempt:2} applied to $y = t(x)$, we have $h(t(t(x))) = h(t(y)) = y = t(x)$
and so $t(x) \in \fix{\concat{t}{h}}$.

$\ref{prelim:prop:compositionidempt:4} \Rightarrow \ref{prelim:prop:compositionidempt:1}$:
We have $h(t(x)) = h(( \move{h}\cap t(X) ) \cup ( \fix{h}\cap t(X) ) )
= h(\move{h} \cap t(X)) \cup h(\fix{h}\cap t(X)) 
= h(\move{h} \cap t(X)) \cup (\fix{h}\cap t(X)) \subseteq \fix{\concat{t}{h}}$.
Thus, by \cref{prelim:prop:idempotent},
$\concat{t}{h}$ is idempotent.
\end{proof}

\begin{corollary} \label{prelim:prop:compositionidempth} \label{prelim:prop:compositionidemptt}
	Let $t, h \colon X \to X$.
	\begin{enumerate}
	\item \label{prelim:prop:compositionidempth:1} For idempotent~$h$, $\concat{t}{h}$ is idempotent
	if and only if $\move{t} \cap h(t(X)) \subseteq \fix{ht}$.
	
	\item \label{prelim:prop:compositionidemptt:1} For idempotent~$t$, $\concat{t}{h}$ is idempotent
	if and only if $h(\move{h} \cap t(X)) \subseteq \fix{\concat{t}{h}}$.
	\end{enumerate}
\end{corollary}
\begin{proof}
\ref{prelim:prop:compositionidempth:1}: 
 For idempotent~$h$, we have $h(t(X)) \subseteq h(X) \subseteq \fix{h}$
 by \cref{prelim:prop:idempotent}.
 Thus, $\fix{t} \cap h(t(X)) \subseteq \fix{h}$ is always fulfilled
 and the claim follows from \cref{prelim:prop:compositionidempt:3} of \cref{prelim:prop:compositionidempt}.

\ref{prelim:prop:compositionidemptt:1}: By \cref{prelim:prop:idempotent}
we have $t(X) \subseteq \fix{t}$.
Since $\fix{h}\cap \fix{t}\subseteq \fix{\concat{t}{h}}$ by \cref{prelim:trans:eqn:fix-comp},
the claim follows from \cref{prelim:prop:compositionidempt:4} of \cref{prelim:prop:compositionidempt}.
\end{proof}

%\begin{remark}
%By the previous corollary, for idempotent~$h$,
%$\concat{t}{h}$ is idempotent if $h(t(X)) \subseteq \fix{t}$
%since then $\move{t} \cap h(t(X)) = \emptyset$.
%But this also follows directly from $h(t(X)) \subseteq \fix{h} \cap \fix{t} \subseteq \fix{\concat{t}{h}}$
%with \crefitem{prelim:prop:idempotent}{prop:idemp:item:img}.
%We note that this condition is not necessary.
%For example, if $h(x) = c$ for all $x \in X$ and $t(c) \ne c$, then
%$h$ and $\concat{t}{h}$ are idempotent and $h(t(X)) \cap \fix{t} = \emptyset$.
%\end{remark}

\begin{corollary} \label{prelim:cor:comp-idempotents-idempotent}
	Let $t, h \colon X \to X$ be idempotent. 
	Then $\concat{t}{h}$ is idempotent if and only if 
	$\move{t} \cap h(\move{h} \cap t(X)) \subseteq \fix{\concat{t}{h}}$.\footnote{Equivalently,
	for $x \in X$, $h(t(x)) \ne t(x)$ and
	$t(h(t(x))) \ne h(t(x))$ imply $h(t(x)) = h(t(h(t(x))))$.}
\end{corollary}
\begin{proof}
	First, suppose $\concat{t}{h}$ is idempotent.
	Since~$h$ and~$t$ are idempotent, by \cref{prelim:prop:compositionidempth} 
	we find $\move{t} \cap h(\move{h} \cap t(X)) \subseteq \fix{\concat{t}{h}}$.
	
	Conversely, suppose $\move{t} \cap h(\move{h} \cap t(X)) \subseteq \fix{\concat{t}{h}}$.
	Since $t$ is idempotent,
	we have $h(\fix{h}\cap t(X)) = \fix{h}\cap t(X) = \fix{h}\cap \fix{t} \subseteq \fix{\concat{t}{h}}$.
	Since~$h$ is idempotent,
	we have $\fix{t} \cap h(t(X)) \subseteq \fix{t}\cap \fix{h} \subseteq \fix{\concat{t}{h}}$.
	We have
	\begin{multline*}
	h(t(X)) = ( \move{t} \cap h(\move{h} \cap t(X)) ) \cup 
	\\ ( \move{t} \cap h(\fix{h}\cap t(X))) \cup ( \fix{t} \cap h(t(X)) )
	\end{multline*}
	Thus, $\concat{t}{h}\subseteq \fix{\concat{t}{h}}$
	and so $\concat{t}{h}$ is idempotent.	
\end{proof}

\begin{corollary} \label{prelim:cor:comp-idempotents-idempotent2}
	Let $t, h \colon X \to X$ be idempotent.
	Then $\concat{t}{h}$ is idempotent if one of the following is true:
\begin{enumerate}
\item \label{prelim:cor:comp-idempotents-idempotent2:1}
$\dupl{h} \subseteq  \fix{\concat{t}{h}}$, i.e.,  $h(t(y)) = y$ for all $y \in \dupl{h}$,

\item \label{prelim:cor:comp-idempotents-idempotent2:2}
 $\excl{t} \subseteq  \fix{\concat{t}{h}}$,

\item \label{prelim:cor:comp-idempotents-idempotent2:3} $ht = th$.
\end{enumerate}
\end{corollary}
\begin{proof}
\ref{prelim:cor:comp-idempotents-idempotent2:1}: 
 Since $h$ is idempotent, we have $h(\move{h} \cap t(X)) \subseteq \dupl{h}$.
 Thus, the claim follows from \cref{prelim:prop:compositionidemptt:1} of \cref{prelim:prop:compositionidemptt}.
 
\ref{prelim:cor:comp-idempotents-idempotent2:2}: With \cref{prelim:prop:idempotent}, 
 we find $\move{t} = \excl{t} \subseteq \fix{\concat{t}{h}}$.
 Then \cref{prelim:cor:comp-idempotents-idempotent}
 implies idempotency of~$\concat{t}{h}$.
 
\ref{prelim:cor:comp-idempotents-idempotent2:3}: We have $h(t(h(t(x)))) = h(h(t(t(x)))) = h(t(x))$ for $x \in X$. 	
\end{proof}

\begin{proposition} \label{prelim:trans:idemp:prop:tX-in-fixh}
	Let $t, h \colon X \to X$ be idempotent.
	Then the following are equivalent.
	\begin{enumerate}
		\item  \label{prelim:trans:idemp:prop:tX-in-fixh:1} $\concat{t}{h} = t$.
		
		\item  \label{prelim:trans:idemp:prop:tX-in-fixh:2} $t(X) \subseteq \fix{h}$.
		
		\item  \label{prelim:trans:idemp:prop:tX-in-fixh:3} $t(X) \cap h^{-1}(\dupl{h}) \subseteq \dupl{h}$.
	\end{enumerate}
\end{proposition}
\begin{proof} 
	$\ref{prelim:trans:idemp:prop:tX-in-fixh:1} \Rightarrow \ref{prelim:trans:idemp:prop:tX-in-fixh:2}$: By assumption,
	we have $h(t(x)) = t(x)$ for all $x \in X$ and so $t(X) \subseteq \fix{h}$.
	
	$\ref{prelim:trans:idemp:prop:tX-in-fixh:2} \Rightarrow \ref{prelim:trans:idemp:prop:tX-in-fixh:3}$: 	
	Let $t(x) \in h^{-1}(\dupl{h})$.
	Since $t(X) \subseteq \fix{h}$, we have $t(x) = h(t(x)) \in \dupl{h}$.
	
	$\ref{prelim:trans:idemp:prop:tX-in-fixh:3} \Rightarrow \ref{prelim:trans:idemp:prop:tX-in-fixh:1}$: 
	Suppose $h(t(x)) \ne t(x)$ for $x \in X$.
	Then $h(t(x)) \in \dupl{h}$ since~$h$ is idempotent,
	but $t(x) \notin \dupl{h}$ since $h$ fixes every point from $\dupl{h}$.
	Thus \cref{prelim:trans:idemp:prop:tX-in-fixh:3} fails.
\end{proof}

Next, we refine \cref{prelim:func:prop:inv-dupl}
for the case that one transformation is idempotent. 
We need the following equation for idempotent $t \colon X \to X$ and $A \subseteq X$ that is easy to verify
\begin{equation} \label{prelim:trans:idemp:eqn:inv}
	t^{-1}(A) = ( A \cap t(X) ) \cup t^{-1}( A \cap \dupl{t} )
\end{equation}

\begin{proposition}  \label{prelim:trans:idemp:prop:inv-dupl}
	Let $t \colon X \to X$ be idempotent and $h \colon X \to Y$.
	Then we have
	\begin{align*}
	(\concat{t}{h})^{-1}(\dupl{\concat{t}{h}}) 
	& = t^{-1}(\dupl{t}) \cup ( h^{-1}( \dupl{h_{|t(X)}} ) \cap t(X) ) \\
	& = t^{-1}(\dupl{t}) \cup h^{-1}( \dupl{\concat{t}{h}})	
	\end{align*}
\end{proposition}
\begin{proof}
	Set $A = h^{-1}(\dupl{h_{|t(X)}})$. 
	Now, we have the following:
	\begin{align*}
		(\concat{t}{h})^{-1}(\dupl{\concat{t}{h}}) 
		& = t^{-1}(\dupl{t}) \cup t^{-1}(A) \tag{Prop.~\ref{prelim:func:prop:inv-dupl}} \\
		& = t^{-1}(\dupl{t}) \cup t^{-1}(A \cap \dupl{t} ) \cup ( A  \cap t(X) ) \tag{Eqn.~\ref{prelim:trans:idemp:eqn:inv}} \\
		& = t^{-1}(\dupl{t}) \cup ( A \cap t(X) )
	\end{align*}
	
	Next, set $B = A \cap t(X) = h_{|t(X)}^{-1}(\dupl{h_{|t(X)}})$.
	Using $h(\dupl{t}) = h_{|t(X)}(\dupl{t})$, \cref{prelim:eqn:finvfinclusion},
	and the idempotency of~$t$ we have the following:
	\begin{align*}
		& t^{-1}(\dupl{t}) \cup h^{-1}(\dupl{\concat{t}{h}}) \\
		& = t^{-1}(\dupl{t}) \cup ( t(X) \cap h^{-1}(\dupl{\concat{t}{h}}) ) \tag{$\excl{t}\subseteq t^{-1}(\dupl{t})$} \\
		& = t^{-1}(\dupl{t}) \cup h_{|t(X)}^{-1}(\dupl{\concat{t}{h}})   \\
		& = t^{-1}(\dupl{t}) \cup h_{|t(X)}^{-1}(h(\dupl{t} \cup \dupl{h_{|t(X)}})) \tag{Eqn.~\eqref{prelim:eqn:dupl0}} \\
		& = t^{-1}(\dupl{t}) \cup h_{|t(X)}^{-1}(h(\dupl{t})) \cup B \\
		& = t^{-1}(\dupl{t}) \cup h_{|t(X)}^{-1}(h_{|t(X)}(\dupl{t})) \cup B \\
		& = t^{-1}(\dupl{t}) \cup \dupl{t} \cup h_{|t(X)}^{-1}(h_{|t(X)}(\dupl{t}) \cap \dupl{h_{|t(X)}}) \cup B \\
		& = t^{-1}(\dupl{t}) \cup \dupl{t} \cup B  \\
		& = t^{-1}(\dupl{t}) \cup B \tag{$\dupl{t}\subseteq t^{-1}(\dupl{t})$}
	\end{align*}
	This shows the claimed identities.	
\end{proof}

\begin{remark}	
	In \cref{prelim:trans:idemp:prop:inv-dupl}, the set
	$h^{-1}(\dupl{h_{|t(X)}}) \cap t(X)$
	may be replaced with either
	$h^{-1}(\dupl{h_{|t(X)}})$
	or
	$h^{-1}(\dupl{\concat{t}{h}}) \cap t(X)$,
	since both contain
	$h^{-1}(\dupl{h_{|t(X)}}) \cap t(X)$
	and are contained in
	$h^{-1}(\dupl{\concat{t}{h}})$.
	However,
	$h^{-1}(\dupl{h_{|t(X)}})$
	cannot be replaced with
	$h^{-1}(\dupl{h})$, see
	\cref{prelim:func:ex:inv-dupl}.
\end{remark}
\end{comment}